\chapter{Chapter05 — Nearly Constant Level One}
%%% generated by the blueprint pipeline from FABL.Chapter05.NearlyConstantLevelOne
\section{Overview}

Declarations of the Lean module \texttt{FABL.Chapter05.NearlyConstantLevelOne}.

\section{Declarations}

\begin{definition}[Minority indicator for a Boolean function]
\lean{FABL.nearlyConstantMinorityIndicator}
\label{FABL.nearlyConstantMinorityIndicator}
\leanok
\uses{FABL.BooleanFunction, FABL.BooleanFunction.toReal, FABL.mean}
For a Boolean function $f$ on the sign cube, define its minority indicator at $x$ to be
$(1-f(x))/2$ when the mean of $f$ is nonnegative, and $(1+f(x))/2$ when that mean is
negative.
\end{definition}

\begin{theorem}[Range of the minority indicator]
\lean{FABL.nearlyConstantMinorityIndicator_mem_Icc}
\label{FABL.nearlyConstantMinorityIndicator_mem_Icc}
\leanok
\uses{FABL.BooleanFunction, FABL.BooleanFunction.toReal, FABL.mean, FABL.nearlyConstantMinorityIndicator}
\difficulty{4}
Let $n$ be a natural number, let $f:\{-1,1\}^n\to\{-1,1\}$ be a Boolean function, and
let $x\in\{-1,1\}^n$. The minority indicator of $f$ at $x$, defined using the sign of
the mean of $f$, lies in the interval $[0,1]$.
\end{theorem}

\begin{theorem}[Mean bound for a Boolean sign function]
\lean{FABL.abs_mean_toReal_le_one}
\label{FABL.abs_mean_toReal_le_one}
\leanok
\uses{FABL.BooleanFunction, FABL.BooleanFunction.toReal, FABL.mean}
\difficulty{4}
Let $n\in\mathbb N$ and let $f\colon\{-1,1\}^n\to\{-1,1\}$. Its uniform mean satisfies
$|2^{-n}\sum_{x\in\{-1,1\}^n}f(x)|\leq 1$.
\end{theorem}

\begin{theorem}[Mean of the minority indicator]
\lean{FABL.mean_nearlyConstantMinorityIndicator}
\label{FABL.mean_nearlyConstantMinorityIndicator}
\leanok
\uses{FABL.BooleanFunction, FABL.BooleanFunction.toReal, FABL.mean, FABL.nearlyConstantMinorityIndicator}
\difficulty{5}
Let $n$ be a nonnegative integer and let $f:\{−1,1\}^n\to\{−1,1\}$. The uniform mean of
its nearly constant minority indicator is $\bigl(1-|\mathbb E f|\bigr)/2$, where
$\mathbb E f$ is the uniform mean of $f$ viewed as real-valued.
\end{theorem}

\begin{theorem}[Nonnegative minority-indicator mean]
\lean{FABL.mean_nearlyConstantMinorityIndicator_nonneg}
\label{FABL.mean_nearlyConstantMinorityIndicator_nonneg}
\leanok
\uses{FABL.BooleanFunction, FABL.abs_mean_toReal_le_one, FABL.mean, FABL.mean_nearlyConstantMinorityIndicator, FABL.nearlyConstantMinorityIndicator}
\difficulty{2}
Let $f$ be a function from the $n$-dimensional sign cube to $-1,1$, and regard its
values as real numbers. Define its minority indicator at each point to be $(1-f(x))/2$
when the mean of $f$ is nonnegative, and $(1+f(x))/2$ otherwise. The mean of this
indicator is nonnegative.
\end{theorem}

\begin{theorem}[Mean bound for the minority indicator]
\lean{FABL.mean_nearlyConstantMinorityIndicator_le}
\label{FABL.mean_nearlyConstantMinorityIndicator_le}
\leanok
\uses{FABL.BooleanFunction, FABL.BooleanFunction.toReal, FABL.mean, FABL.mean_nearlyConstantMinorityIndicator, FABL.nearlyConstantMinorityIndicator}
\difficulty{2}
Let $n$ be a nonnegative integer, let $f$ be a function from the $n$-dimensional sign
cube $\\{-1,1\}^n$ to $\{-1,1\}$, and let $\delta\in\mathbb R$. If $1-\delta\le |\mathbb
E f|$, where the mean is taken uniformly over the sign cube, then the uniform mean of
the minority indicator of $f$ is at most $\delta/2$.
\end{theorem}

\begin{theorem}[Singleton coefficient of the minority indicator]
\lean{FABL.fourierCoeff_nearlyConstantMinorityIndicator_singleton}
\label{FABL.fourierCoeff_nearlyConstantMinorityIndicator_singleton}
\leanok
\uses{FABL.BooleanFunction, FABL.BooleanFunction.toReal, FABL.expect_monomial, FABL.fourierCoeff, FABL.mean, FABL.monomial, FABL.nearlyConstantMinorityIndicator}
\difficulty{5}
Let $n$ be a nonnegative integer, let $f:\{-1,1\}^n\to\{\pm1\}$ be a Boolean function,
and let $i$ be a coordinate. The Fourier coefficient of its minority indicator at
$\{i\}$ is $-\widehat f(\{i\})/2$ if the mean of the real-valued form of $f$ is
nonnegative, and is $\widehat f(\{i\})/2$ otherwise.
\end{theorem}

\begin{theorem}[Level-one weight of the minority indicator]
\lean{FABL.fourierWeightAtLevel_one_eq_four_mul_nearlyConstantMinorityIndicator}
\label{FABL.fourierWeightAtLevel_one_eq_four_mul_nearlyConstantMinorityIndicator}
\leanok
\uses{FABL.BooleanFunction, FABL.BooleanFunction.toReal, FABL.fourierCoeff_nearlyConstantMinorityIndicator_singleton, FABL.fourierWeightAtLevel, FABL.fourierWeightAtLevel_one_eq_sum_singleton, FABL.mean, FABL.nearlyConstantMinorityIndicator}
\difficulty{4}
Let $f:\{-1,1\}^n\to\{-1,1\}$, and define $h(x)$ to be $(1-f(x))/2$ if the uniform mean
of $f$ is nonnegative, and $(1+f(x))/2$ otherwise. The level-one Fourier weight of $f$
is four times the level-one Fourier weight of $h$, where level-one weight is the sum of
squared Fourier coefficients over singleton coordinate sets.
\end{theorem}

\begin{theorem}[Monotonicity of a logarithmic expression]
\lean{FABL.sq_mul_log_inv_mono}
\label{FABL.sq_mul_log_inv_mono}
\leanok
\difficulty{5}
For real numbers $x$ and $y$ with $0<x$, $x\le y$, and $y\le \tfrac12$, one has
$x^2\log(1/x)\le y^2\log(1/y)$.
\end{theorem}

\begin{theorem}[Fourier weight at level one le of abs mean ge]
\lean{FABL.fourierWeightAtLevel_one_le_of_abs_mean_ge}
\statementsource{Analysis-of-Boolean-Functions-by-Ryan-ODonnell}{
  pdf-ffbe9c7b3306-p126-b002,
  pdf-ffbe9c7b3306-p126-b003,
  pdf-ffbe9c7b3306-p128-b004
}
\label{FABL.fourierWeightAtLevel_one_le_of_abs_mean_ge}
\leanok
\uses{FABL.BooleanFunction, FABL.BooleanFunction.toReal, FABL.abs_mean_toReal_le_one, FABL.fourierWeightAtLevel, FABL.fourierWeightAtLevel_one_eq_four_mul_nearlyConstantMinorityIndicator, FABL.mean, FABL.mean_nearlyConstantMinorityIndicator, FABL.mean_nearlyConstantMinorityIndicator_le, FABL.mean_nearlyConstantMinorityIndicator_nonneg, FABL.nearlyConstantMinorityIndicator, FABL.nearlyConstantMinorityIndicator_mem_Icc, FABL.sharpLevelOneInequality, FABL.sharpLevelOneInequality_eq_zero, FABL.sq_mul_log_inv_mono}
\difficulty{6}
% dataset-pipeline: stub -- prose pending build_informal_statements.py
Statement pending regeneration.
\end{theorem}
